\begin{tabular}{>{\raggedright\arraybackslash}p{0.22\textwidth}>{\raggedright\arraybackslash}p{0.62\textwidth}}
\arrayrulecolor{gray!60}\hline
\rowcolor{grammarheadcolor}$\odot$ & $\hat{\odot}(v, v')$ \\
\pyinline{==} & $\val{\pyinline{False}}$ if both operands are \pyinline{nan}, and $\eqOp{\Sigma}(v, v')$ otherwise \\
\pyinline{!=} & $\val{\pyinline{True}}$ if both operands are \pyinline{nan} or $\eqOp{\Sigma}(v, v') = \val{\pyinline{False}}$, $\val{\pyinline{False}}$ if $\eqOp{\Sigma}(v, v') = \val{\pyinline{True}}$, and otherwise $\eqOp{\Sigma}(v, v')$ \\
\pyinline{in} & $\contains{\Sigma}(v', v)$ \\
\pyinline{not in} & $\val{\pyinline{True}}$ if $\contains{\Sigma}(v', v) = \val{\pyinline{False}}$, $\val{\pyinline{False}}$ if $\contains{\Sigma}(v', v) = \val{\pyinline{True}}$, and otherwise $\contains{\Sigma}(v', v)$ \\
\pyinline{<}\quad\pyinline{<=}\quad\pyinline{>}\quad\pyinline{>=} & \todo{Ordering, as in Python\added{, except $\aborts{\errExcluded}$ where an operand is a Boolean, which Python treats as the integer 0 or 1}} \\
\pyinline{+}\quad\pyinline{-}\quad\pyinline{*}\quad\pyinline{/}\quad\pyinline{//}\quad\pyinline{\%}\quad\pyinline{**} & \todo{Arithmetic, as in Python. Float precision is left to the implementation. $\aborts{\errZeroDiv}$ where the divisor is zero, \added{$\aborts{\errOverflow}$ where Python raises the exception of that name, such as for an integer too large to convert to a float, $\aborts{\errExcluded}$ where an operand is a Boolean, which Python treats as the integer 0 or 1,} and $\aborts{\errType}$ where an operand is not a number}
\end{tabular}
